\renewcommand{\theenumi}{\thesubsection.\arabic{enumi}}
\renewcommand{\labelenumi}{\textbf{\theenumi}}
\begin{enumerate}[leftmargin=0pt,itemsep=6pt]
\item
\paragraph*{Nonlinear parameter identification and root multiplicity}
For $P_p(x)=px^2-2\sqrt5\,px+15$, the repeated-root constraint is
\begin{equation}
F(p)=\Delta(p)=20p^2-60p=0,\qquad
\Delta=b^2-4ac.
\label{eq:quad-general}
\end{equation}
Newton--Raphson on the parameter equation uses
\begin{equation}
p_{n+1}=p_n-\frac{20p_n(p_n-3)}{40p_n-60}.
\label{eq:newton-multiplicity-01}
\end{equation}
For a double root $x_\star$ of $g(x)=c(x-x_\star)^2$, Newton's error obeys
$e_{n+1}=e_n/2$, so convergence is linear rather than quadratic.
\item
\begin{equation}
\tan\theta=\frac{\text{opposite}}{\text{adjacent}},\quad
\sin\theta=\frac{\text{opposite}}{\text{hypotenuse}},\quad
\cos\theta=\frac{\text{adjacent}}{\text{hypotenuse}}.
\label{eq:trig-ratios}
\end{equation}
Useful values are $\tan30^\circ=1/\sqrt3$, $\tan60^\circ=\sqrt3$, $\cot30^\circ=\sqrt3$, $\cot60^\circ=1/\sqrt3$, $\sin30^\circ=1/2$, and $\cos60^\circ=1/2$. For a right triangle,
\[c^2=a^2+b^2,\qquad \|(x,y)^T\|=\sqrt{x^2+y^2}.\]

\paragraph*{Probability}
For equally likely outcomes,

\item
\paragraph*{Push-forward product distribution}
For independent $X,Y\sim\operatorname{Unif}\{1,\ldots,6\}$ and $Z=XY$,
\begin{align}
P_Z(m)&=\frac1{36}\sum_{i=1}^{6}\sum_{j=1}^{6}\mathbf1_{\{ij=m\}},\\
G_Z(w)&=E[w^Z]=\frac1{36}\sum_{i=1}^{6}\sum_{j=1}^{6}w^{ij}.
\label{eq:product-pgf}
\end{align}
The coefficient $[w^m]G_Z(w)$ is $P(Z=m)$, and $G_Z(1)=1$ verifies total probability.
\item
\begin{equation}
\theta_{\mathrm{tangent-chord}}=\theta_{\mathrm{alternate\ segment}}.\label{eq:tangent-chord}
\end{equation}

For a triangle with inradius $r$ and semiperimeter $s$,
\addtocounter{enumi}{1}
\item
\begin{equation}
[ABC]=rs,\qquad s=\frac{AB+BC+CA}{2}.
\label{eq:triangle-inradius}
\end{equation}
Tangents from the same external point have equal lengths:
\begin{equation}
PA=PB.\label{eq:equal-tangents}
\end{equation}

\paragraph*{Arithmetic progressions}
For first term $a$ and common difference $d$,

\addtocounter{enumi}{1}
\item
\paragraph*{Jacobi iteration and convergence}
For $Ax=b$, $A=D+R$ with diagonal $D$, Jacobi iteration is
\begin{equation}
x^{(k+1)}=-D^{-1}Rx^{(k)}+D^{-1}b,\qquad
\rho(-D^{-1}R)<1\Longrightarrow x^{(k)}\to A^{-1}b.
\label{eq:jacobi-08}
\end{equation}
The residual $r^{(k)}=b-Ax^{(k)}$ and the increment
$\|x^{(k+1)}-x^{(k)}\|_\infty$ provide numerical stopping checks.
\item
\begin{equation}
\overrightarrow{BA}\cdot\overrightarrow{BC}=0.\label{eq:right-angle-dot}
\end{equation}

For $A(x_1,y_1)$, $B(x_2,y_2)$, and $C(x_3,y_3)$, collinearity is equivalent to

\item
\paragraph*{Rank and collinearity}
For two displacement rows $u,v$ in $\mathbb R^2$, the coordinate matrix
\[
M=\begin{pmatrix}u_1&u_2\\v_1&v_2\end{pmatrix}
\]
has rank one exactly when the nonzero row vectors are linearly dependent:
\begin{equation}
\operatorname{rank}(M)=1\iff\det M=0.
\label{eq:collinear-det}
\end{equation}
\item
\paragraph*{Unilateral $Z$-transform of an arithmetic progression}
For the causal sequence $x[n]=a+(n-1)d$, $n\ge1$,
\begin{align}
X(z)&=\mathcal Z_u\{x[n]\}
=\frac{az^{-1}}{1-z^{-1}}+\frac{dz^{-2}}{(1-z^{-1})^2},
\label{eq:ap-question-11}\\
S(z)&=\frac{X(z)}{1-z^{-1}}
\end{align}
is the transform of its causal accumulated sum.
\addtocounter{enumi}{1}
\item
\begin{equation}
\text{distance}=\text{speed}\times\text{time},\qquad
\text{average speed}=\frac{\text{total distance}}{\text{total time}}.
\label{eq:speed-time}
\end{equation}

\paragraph*{Coordinate geometry as linear algebra}
Represent $P(x,y)$ by $\vec p=\myvec{x\\y}$. Then

\item
\begin{equation}
\overrightarrow{PQ}=\vec q-\vec p,\qquad
\langle\vec u,\vec v\rangle=u_xv_x+u_yv_y.
\label{eq:vector-basics}
\end{equation}

If $P$ divides $AB$ internally in the ratio $AP:PB=m:n$, then

\begin{equation}
\vec p=\frac{n\vec a+m\vec b}{m+n}=\vec a+\frac{m}{m+n}(\vec b-\vec a).
\label{eq:section-vector}
\end{equation}

\paragraph*{Probability}

If $E_1,E_2,E_3$ form an exhaustive set of mutually exclusive events, then
\item
\paragraph*{Inverse probability and integer feasibility}
For exhaustive, disjoint colour classes,
\begin{equation}
P(R)+P(B)+P(O)=1,\qquad
R=\frac N4,\quad B=\frac N3,\quad O=N-R-B.
\label{eq:three-events}
\end{equation}
Integer counts require $N$ to satisfy all divisibility and nonnegativity constraints.
\item
\begin{equation}
A_{\mathrm{sector}}=\frac{\theta}{360^\circ}\pi r^2,\qquad
A_{\mathrm{segment}}=A_{\mathrm{sector}}-A_{\mathrm{triangle}}.
\label{eq:sector-segment}
\end{equation}

The curved surface areas of a cylinder and a cone are

\item
\begin{equation}
A_{\mathrm{cyl}}=2\pi rh,\qquad A_{\mathrm{cone}}=\pi r\ell,\qquad \ell=\sqrt{r^2+h^2}.
\label{eq:cone-formula}
\end{equation}

\addtocounter{enumi}{4}
\item
\paragraph*{Finite-support discrete-time sequence}
For $x[k]=a+kd$ on $0\le k\le N-1$, its finite $Z$-transform is
\[
\begin{align}
X(z)&=a\frac{1-z^{-N}}{1-z^{-1}}
+d\frac{z^{-1}\left(1-Nz^{-(N-1)}+(N-1)z^{-N}\right)}{(1-z^{-1})^2},
\label{eq:finite-z-22}\\
S(z)&=\frac{X(z)}{1-z^{-1}}.
\end{align}
The discrete-time accumulator sums the finite-support record without adding samples beyond its terminal index.
\addtocounter{enumi}{4}
\item
\begin{equation}
x=h\cot\theta.\label{eq:horizontal-range}
\end{equation}

\paragraph*{Probability}

For events $A$ and $B$,

\item
\paragraph*{Indicator variables and inclusion--exclusion}
For event indicators,
\begin{align}
I_{A\cup B}&=I_A+I_B-I_AI_B,\\
P(A\cup B)&=E[I_A]+E[I_B]-E[I_AI_B].
\label{eq:prob-union}
\end{align}
For disjoint events the product term is zero; a complement indicator is $1-I_A$.
\item
\paragraph*{Determinant area as a nonlinear parameter constraint}
For vertices $P_1,P_2,P_3\in\mathbb R^2$,
\begin{equation}
2A=\left|\det\begin{pmatrix}
x_2-x_1&y_2-y_1\\
x_3-x_1&y_3-y_1
\end{pmatrix}\right|.
\label{eq:triangle-det}
\end{equation}
When the coordinates depend on a parameter, solve the resulting determinant
equation and check that each candidate gives a nonzero determinant.
\item
\begin{equation}
A_{\mathrm{shaded}}=A_{\mathrm{outer}}-A_{\mathrm{inner}}.\label{eq:shaded-area}
\end{equation}

The corresponding volume formula for a cylinder is

\item
\begin{equation}
V_{\mathrm{cyl}}=\pi r^2h.
\label{eq:cylinder-formula}
\end{equation}

For a hemisphere and sphere,

\begin{equation}
V_{\mathrm{hemisphere}}=\frac23\pi r^3,\qquad V_{\mathrm{sphere}}=\frac43\pi r^3,
\qquad A_{\mathrm{sphere}}=4\pi r^2.
\label{eq:sphere-formula}
\end{equation}

For a cylinder with two hemispherical cavities of radius $r$ scooped from its ends,

\begin{equation}
V_{\mathrm{remaining}}=\pi R^2h-\frac43\pi r^3.
\label{eq:scooped-cylinder}
\end{equation}
\end{enumerate}
